\section{Finite filtered algebras: a common lifting theorem}
\label{sec:filtered-lifting}

The existing Lagrange-inversion packet constructs coefficient prefixes by
fixed-point iteration, and the Puiseux packet reduces branches to simple-root
equations before applying a coefficient recurrence.  The Bell packet uses
finite exponential-series identities.  These are substantial existing
arguments; the purpose of this section is to expose a common reusable
arithmetic theorem, including simultaneous polynomial systems and modular
lifting \cite{repoLagrange,repoPuiseux,repoBell}. The classical Hensel principle is not new. The contribution here
is its organization as one division-free finite program, one bounded
invariant, and one coherence argument suitable for iPRA.

\subsection{The exact arithmetic interface}

Fix a primitive-recursively presented commutative ring $B$.  Its element
codes, validity predicate, equality test, and ring operations are fixed
primitive recursive data, and the ring and congruence laws have supplied
iPRA derivations.  The displayed equalities below mean equality in this
presentation; a canonical representation is convenient but unnecessary.
Let $N\ge1$ be a numerical precision parameter.  Assume that
$F^kB$ is a uniformly primitive recursive membership test, with proofs of
the following laws:
\begin{enumerate}
\item $F^0B=B$, $F^{k+1}B\subseteq F^kB$, and $F^NB=\{0\}$;
\item every $F^kB$ is an additive subgroup;
\item $uv\in F^{i+j}B$ whenever $u\in F^iB$ and $v\in F^jB$.
\end{enumerate}
One can set $F^kB=\{0\}$ for $k\ge N$.  Because $F^0B=B$, the last
condition also gives closure under multiplication by an arbitrary ring
element.  Matrices and vectors belong to a filtration level when each of
their finitely many entries does.

These are \emph{supplied numerical contracts}.  This definition does not
assert that membership in powers of an arbitrary ideal is decidable.
The ring and filtration may depend on numerical parameters, including
$N$, provided all operations and membership tests are jointly primitive
recursive and the same displayed laws are proved uniformly.  The concrete
jet and residue-ring examples below have this property.

A polynomial in $m$ variables is a finite list of terms
$(c,e_1,\ldots,e_m)$, interpreted as
$cX_1^{e_1}\cdots X_m^{e_m}$.  Missing coefficients are zero; duplicate
monomials are allowed.  Evaluation, formal partial derivatives, vector
evaluation, and finite matrix multiplication are primitive recursive
bounded folds.  No general-purpose primitive recursive program evaluator
is used.  The polynomials themselves, their dimensions, and their degrees
may be variable numerical inputs.

\begin{lemma}[Finite polynomial difference and Taylor identities]
\label{lem:filtered-taylor}
Let $P\in B[X_1,\ldots,X_m]^m$, and let $J(x)$ denote its Jacobian.
If $h\in(F^kB)^m$, then
\[
 P(x+h)-P(x)-J(x)h\in(F^{2k}B)^m.                 \tag{T}
\]
If $x-a\in(F^1B)^m$, then every entry of $J(x)-J(a)$ belongs to $F^1B$.
Both claims have uniform iPRA derivations from the presentation contracts.
\end{lemma}

\begin{proof}
Expand one monomial
$c\prod_j(x_j+h_j)^{e_j}$ by the finite binomial theorem.  The term using
no $h_j$ is its value at $x$.  The terms using exactly one $h_j$ have sum
equal to its formal derivative applied to $h$.  Every remaining term has
at least two factors among the $h_j$, counted with multiplicity, and
therefore belongs to $F^{2k}B$.  Multiplication by the other factors
preserves membership, and summing finitely many such terms does so too.
Adding the finitely many monomials proves (T).

For the second assertion apply the elementary finite identity
$u^r-v^r=(u-v)\sum_{j<r}u^jv^{r-1-j}$ to one variable at a time in
each monomial of a partial derivative.  The resulting difference is a
finite sum of multiples of the coordinates $x_j-a_j$, hence belongs to
$F^1B$.  The proof uses no division and is valid in every characteristic.

The binomial and telescoping identities are proved by induction on natural
exponents.  Their induction formulas are equalities of coded ring elements.
The further invariants assert that all terms in a finite accumulated list
have the stated membership; they are bounded formulas and have Boolean
quantifier-free forms.  Thus the finite expansion argument gives actual
iPRA proof instructions uniformly in the polynomial code, rather than
merely an appeal to a general algebra theorem in the ambient metatheory.
\end{proof}

\subsection{A frozen inverse is enough}

\begin{theorem}[Finite filtered lifting]
\label{thm:filtered-lifting}
Let $a\in B^m$ and $U\in B^{m\times m}$ satisfy
\[
 P(a)\in(F^1B)^m,\qquad
 I-J(a)U\in(F^1B)^{m\times m},\qquad
 I-UJ(a)\in(F^1B)^{m\times m}.                 \tag{H}
\]
Define the primitive recursive sequence
\[
 b_0=a,\qquad b_{r+1}=b_r-U P(b_r).            \tag{L}
\]
Then $b_{N-1}$ is a root of $P$, and is the unique root $x$ satisfying
$x-a\in(F^1B)^m$.  No derivative inverse is recomputed during the
iteration, and no division occurs after the data $a,U$ are supplied.

More precisely, for every $r<N$,
\[
 b_r-a\in(F^1B)^m,\qquad
 P(b_r)\in(F^{r+1}B)^m,\qquad
 b_{r+1}-b_r\in(F^{r+1}B)^m.                 \tag{I}
\]
The first two assertions suffice at $r=N-1$; the last assertion is also
meaningful there if one more iterate is formed.
\end{theorem}

\begin{proof}[Existence and the precision invariant]
Put $x=b_r$, $v=P(x)$, and $h=-Uv$.  Suppose $x-a\in F^1$ and
$v\in F^k$, where $k\ge1$.  Matrix multiplication preserves membership
in $F^k$, so $h\in F^k$.  In particular $x+h-a\in F^1$.
Lemma~\ref{lem:filtered-taylor} gives
\[
 P(x+h)=(I-J(a)U)v-(J(x)-J(a))Uv+R,
 \qquad R\in(F^{2k}B)^m.                    \tag{1}
\]
The first term belongs to $F^{k+1}$ by (H).  The second belongs to
$F^{k+1}$ because $J(x)-J(a)\in F^1$.  The remainder does as well:
$2k\ge k+1$ and the filtration is decreasing.  Thus one iteration
raises the certified residual precision by at least one level.

Initially the first two assertions of (I) hold by (H).  The preceding
argument proves the induction step, and the defining equation (L) proves
the increment assertion.  At $r=N-1$ the residual belongs to $F^NB=\{0\}$.
If $N=1$, no update is performed and the hypothesis already says $P(a)=0$.
The proof uses the right inverse congruence $J(a)U\equiv I$; the left
inverse congruence will be used for uniqueness.
\end{proof}

\begin{proof}[Uniqueness]
Let $x,y$ be roots with $x-a,y-a\in F^1$.  We show inductively that
$d=x-y\in F^k$ for $1\le k\le N$.  The case $k=1$ is immediate.
If $d\in F^k$, the Taylor identity at $y$ yields
\[
 0=P(x)-P(y)=J(y)d+R,\qquad R\in F^{2k}.
\]
Multiplying by $U$ and rearranging gives
\[
 d=(I-UJ(a))d-U(J(y)-J(a))d-UR.              \tag{2}
\]
Each term on the right lies in $F^{k+1}$, for exactly the three reasons
used in (1).  This proves the step.  At level $N$, $d=0$, hence $x=y$.
The induction is finite and contains no least-counterexample principle,
choice, field cancellation, or existence assertion about a completion.
\end{proof}

The linear precision gain is a deliberate safe bound.  It need not be
quadratic when $U$ is only an inverse modulo $F^1$.  For example, in
$B=\mathbb Z/q^N\mathbb Z$, take $P(X)=X-q$, $a=0$, and $U=1+q$.
The residual satisfies $P(b_r)=(-q)^{r+1}$.  For $r+1<N$ it is generally
nonzero modulo $q^N$, so the bound of $N-1$ updates is attained.

\paragraph{Why this is an iPRA theorem scheme.}
Fix the final precision $N$ before induction.  At stage $r$, use
\[
 r<N\ \longrightarrow\
 \bigl(\operatorname{ValidVec}(b_r)\wedge
       b_r-a\in F^1\wedge P(b_r)\in F^{r+1}\bigr)
                                                               \tag{3}
\]
as the readable invariant.  Its vector and matrix conditions range only
over their coded finite lengths.  Booleanization converts (3) to a
quantifier-free equation.  The program symbols in (L) are obtained by one
primitive recursion; matrix and polynomial operations are previously fixed
PR definitions.  The existence proof consequently compiles by
\texttt{PRA.qfInd}.  For uniqueness, fix $x,y,N$ as parameters and use the
Booleanization of $k\le N\to x-y\in F^k$, treating the step from
$k=0$ separately by the known congruence modulo $F^1$.
The contracts for the ring and
filtration remain explicit hypotheses or previously supplied iPRA lemmas.
The argument does not deduce object-theory provability merely from a
terminating Lean implementation.

\subsection{Coherence without selecting an infinite object}

\begin{corollary}[Transport of the unique lift]
\label{cor:filtered-coherence}
Suppose $\rho:B\to C$ is a fixed PR map between two PR-presented
filtered rings, with supplied iPRA proofs that it is a ring homomorphism,
sends $F^1B$ into $F^1C$, and sends the polynomial system
$P$ to $Q$.  Let $b$ be the root produced in $B$.  Suppose a second valid
lifting instance in $C$ has initial approximation $c$ satisfying
$\rho(a)-c\in F^1C$.  Then $\rho(b)$ equals its produced root.
\end{corollary}

\begin{proof}
Finite polynomial evaluation commutes with the homomorphism, so
$Q(\rho(b))=\rho(P(b))=0$.  Moreover $\rho(b)-c\in F^1C$.
The uniqueness part of Theorem~\ref{thm:filtered-lifting}, applied only in
$C$, identifies this root with the one produced there.  Different lifts
of the initial approximation or of the inverse matrix therefore give the
same answer whenever they satisfy their local hypotheses.
\end{proof}

In a tower of finite quotients, construct the root independently at each
requested precision.  The corollary proves that truncating one answer gives
the answer at lower precision.  This produces a coherent family by one
fixed PR formula.  There is no step of choosing one element from each of
infinitely many nonempty sets.  Inside iPRA, coherence is an arithmetic
statement about any two requested numerical precisions.

\subsection{Ordinary jets and simultaneous implicit functions}

Take
\[
 B_N=R[T]/(T^N),\qquad
 F^kB_N=\{f:[T^j]f=0\text{ for every }j<\min(k,N)\},
                                                               \tag{4}
\]
where $R$ is any fixed PR-presented commutative ring.  An element is a
length-$N$ coefficient list.  Addition is entrywise and multiplication is
the truncated convolution
\[
 [T^n](fg)=\sum_{j=0}^n [T^j]f\,[T^{n-j}]g\quad(n<N).
\]
All filtration laws follow by bounded sums.  For the product law, if
$n<i+j$, every summand has either its first index below $i$ or its second
index below $j$; that factor is zero.  This is a bounded decidable case
split.  Truncation is a ring homomorphism.

\begin{corollary}[Finite simultaneous formal implicit functions]
\label{cor:formal-implicit}
Let $P(T,Y_1,\ldots,Y_m)$ be a polynomial system over $R$.  Supply
$a\in R^m$ and $V\in R^{m\times m}$ such that
\[
 P(0,a)=0,\qquad VJ_Y P(0,a)=J_Y P(0,a)V=I.
\]
For each $N\ge1$, a uniform primitive recursive construction returns the
unique length-$N$ vector of coefficient lists $Y^{[N]}$ with constant
coefficient $a$ and
\[
 P(T,Y^{[N]})=0\pmod {T^N}.
\]
The lists agree under truncation.  The result is valid in arbitrary
characteristic and uses no division by coefficient indices.
\end{corollary}

\begin{proof}
Embed $a,V$ as constant lists and apply Theorem~\ref{thm:filtered-lifting}
in (4).  The hypotheses become its three membership tests at level one.
Apply Corollary~\ref{cor:filtered-coherence} to the truncation maps.
\end{proof}

For a general system $Y=a+T G(T,Y)$, the Jacobian at $T=0$ is the identity.
Thus the theorem handles mutually recursive algebraic coefficient
definitions with no inversion input at all.  The single equation
$Y=1+TY^2$ gives the Catalan recurrence.  The pair
\[
 Y=1+TZ^2,\qquad Z=1+TYZ
\]
illustrates that simultaneous recursion is treated by the same program
and proof; no ordering of the unknown functions has to be chosen.

For inversion, if $f(0)=c$ and $vc=1$, take $P(T,Y)=f(T)Y-1$ and
$a=U=v$.  Put $q=1-vf\in F^1$.  Iteration gives explicitly
\[
 Y^{[N]}=v\sum_{j=0}^{N-1}q^j,
 \qquad fY^{[N]}=(1-q)\sum_{j=0}^{N-1}q^j=1-q^N=1.
\]
This proves the usual inverse-series construction from a finite geometric
identity.  For $P=Y^2-(1+T)$, $a=1$ and $U=1/2$ give the formal square-root
coefficients over a ring where $2$ has a supplied inverse.  Over a ring
where it has none, this particular instance does not satisfy the theorem's
hypotheses; the interface does not hide the obstruction.

These are statements about exact finite jets.  They do not imply an
analytic convergence radius, a real root, or completeness of all singular
branches.  The Puiseux packet's ramification and exact coefficient-field
backend, and its majorant estimates, remain separate mathematics.
Likewise, this theorem does not reprove the Lagrange coefficient formulas;
it supplies a common construction and uniqueness layer beneath them.

\subsection{Modular Hensel lifting, including composite bases}

Let $q\ge2$ and
\[
 B_N=\mathbb Z/q^N\mathbb Z,\qquad
 F^kB_N=\{x:q^{\min(k,N)}\mid x\},            \tag{5}
\]
using the canonical representative $0\le x<q^N$ in the membership test.
Changing a representative by a multiple of $q^N$ preserves this test.
Divisibility, modular operations, powers, and their finite arithmetic
laws are primitive recursive.  The filtration product law follows by
multiplying two divisibility witnesses; if $i+j\ge N$, the product
is zero in $B_N$.

\begin{corollary}[Finite modular lifting]
\label{cor:modular-hensel}
Let $P\in\mathbb Z[Y_1,\ldots,Y_m]^m$ and supply an integer vector $a$
and matrix $U$ satisfying
\[
 P(a)\equiv0\pmod q,\qquad
 J(a)U\equiv UJ(a)\equiv I\pmod q.
\]
For each $N\ge1$, the recurrence
\[
 b_0=a\bmod q^N,\qquad
 b_{r+1}=b_r-UP(b_r)\bmod q^N
\]
returns after $N-1$ updates the unique root modulo $q^N$ in the residue
class of $a$ modulo $q$.  These outputs are compatible under reduction.
No primality assumption on $q$ is required.
\end{corollary}

\begin{proof}
Apply Theorem~\ref{thm:filtered-lifting} in (5), and use the natural
reduction homomorphism for compatibility.  All inverse information is
contained in the displayed finite congruence, which can be checked before
the iteration.  No claim is made that every root modulo $q$ has such a
Jacobian inverse.
\end{proof}

For example, take $q=4$, $P(X)=X^3-X-2$, $a=2$, $U=3$, and $N=3$.
Here $P(2)=4$ and $3P'(2)=33\equiv1\pmod4$.  Modulo $64$ the two
updates are $2\mapsto54\mapsto22$.  Directly,
$22^3-22-2=10624=166\cdot64$.  The theorem proves uniqueness among the
residues congruent to $2$ modulo $4$, not uniqueness among all residues.

\subsection{Hurwitz coefficients and differential equations without division}

There is a second coefficient algebra particularly suitable for integral
combinatorics.  A length-$N$ Hurwitz list has multiplication
\[
 (a\star b)_n=\sum_{j=0}^n\binom nj a_jb_{n-j}\qquad(n<N),       \tag{6}
\]
with identity $(1,0,\ldots,0)$.  Binomial coefficients are defined by
Pascal recursion.  This algebra is valid over any commutative ring and
does not require factorial denominators.  Commutativity follows from
$\binom nj=\binom n{n-j}$.  Associativity follows by collecting the terms
$a_i b_j c_k$ with $i+j+k=n$ and using
\[
 \binom n{i+j}\binom{i+j}i
   =\binom ni\binom{n-i}j.
\]
This identity counts the same ordered division of $n$ positions into
three labeled parts; alternatively it follows by Pascal induction.  Both
the counting sets and the induction invariants are finite.  The filtration
by vanishing initial coefficients satisfies exactly the laws used above,
so Theorem~\ref{thm:filtered-lifting} applies to this algebra too.

Differentiation is the shift $(Da)_n=a_{n+1}$.  Its finite-precision type
must be respected:
\[
 D:\operatorname{Hurwitz}_{N+1}(R)\longrightarrow
       \operatorname{Hurwitz}_{N}(R).                         \tag{7}
\]
It is not an endomorphism of a fixed truncation. Explicitly, with
$a,b$ of length $N+1$, the typed identity is
\[
 D(a\star_{N+1}b)=Da\star_N\pi_N(b)+\pi_N(a)\star_N Db,
\]
where $\pi_N$ truncates to length $N$. Pascal's identity gives
$D(a\star b)=Da\star b+a\star Db$ at every coefficient for which the
needed input coefficients have been supplied.  Indeed, split
$\binom{n+1}j$ into $\binom nj+\binom n{j-1}$ and reindex the two finite
sums.  The demand for one extra coefficient is part of the contract.

\begin{theorem}[Triangular Hurwitz differential systems]
\label{thm:hurwitz-ode}
Fix a PR-presented commutative ring $R$, an initial vector $c\in R^m$,
and a polynomial system $G(T,Y)$ whose coefficient lists are supplied at
the requested finite precision.  There is a unique coefficient prefix
through degree $N$ satisfying
\[
 Y_0=c,\qquad DY=G(T,Y)\quad\text{through degree }N-1.          \tag{8}
\]
The prefix is produced by a PR recurrence, uniformly in $N$, the finite
polynomial code, and the finite coefficient input.  Compatible input
prefixes give compatible output prefixes.  No integer division is used.
\end{theorem}

Here the time variable $T$ is the Hurwitz list $(0,1,0,\ldots)$.
If the coefficients of $G$ are themselves time-dependent lists, they
are additional finite inputs, with their prefixes supplied through
degree $N-1$. Evaluation uses the Hurwitz product throughout. In
particular $T^{\star r}$ has coefficient $r!$ in position $r$;
one does not identify it with the ordinary monomial basis in positive
characteristic.

\begin{proof}
Assume coefficients $Y_0,\ldots,Y_n$ are already stored.  Form each
coefficient of the finite expression $G(T,Y)$ through degree $n$, using
(6) and the stored lists, and set
\[
 Y_{n+1}=[G(T,Y)]_n.                                           \tag{9}
\]
This expression only reads coefficients through degree $n$: at a product,
every convolution index is at most $n$, and addition and constants have
the same property.  Thus (9) is a triangular recurrence, not an implicit
unbounded search.  A single recursion accumulates the finite vector of
prefixes.  Each step makes one new instance of (8) true and preserves all
earlier instances because the earlier coefficients are unchanged.

For uniqueness, suppose another finite prefix satisfies (8).  Its zeroth
coefficient is $c$.  If its coefficients agree through $n$, finite
evaluation of $G$ agrees at coefficient $n$, so (9) forces agreement at
$n+1$.  Bounded-prefix induction proves equality of the whole requested
prefix.  The same argument proves compatibility upon shortening the
coefficient input.  The invariants are conjunctions over $j\le n$ and
$i<m$ of coded ring equalities, so they have quantifier-free Boolean forms.
\end{proof}

For $Y'=A\star Y$, $Y_0=1$, formula (9) becomes
\[
 Y_{n+1}=\sum_{j=0}^n\binom nj A_jY_{n-j}.                      \tag{10}
\]
Taking $A_j=1$ gives the Bell recurrence, while $Y'=Y\star Y$, $Y_0=1$,
gives $Y_n=n!$: every one of the $n+1$ summands
$\binom nj j!(n-j)!$ equals $n!$.
The simultaneous system $Y'=Z$, $Z'=-Y$ gives the usual sine and cosine
coefficient patterns from their initial values without dividing by any
factorial.

Over a ring containing the rationals, the map
$a\mapsto\sum a_n T^n/n!$ identifies (6) with ordinary multiplication
of exponential generating functions.  Over $\mathbb Z$ or a ring of
positive characteristic, that notation is only a motivation; the
division-free algebra (6) is the definition.  Analytic convergence and
tail estimates are additional contracts.  This recovers the finite
recurrence layer of the Bell packet and generalizes it to arbitrary
polynomial differential systems, without claiming its analytic asymptotic
results anew.

\subsection{Lean definitions and completion obligations}

The finite data layer should define the following objects before invoking
any host-language theorem about rings:
\begin{enumerate}
\item \texttt{RingCode}: numerical validity and equality tests; codes for
zero, one, addition, negation and multiplication; explicit PR syntax for
each of these fixed operations.
\item \texttt{FilteredCode}: a precision parameter, the Boolean predicate
\texttt{member k x}, and the supplied object-iPRA derivations of the
filtration laws above.  \texttt{VecMember} and \texttt{MatMember} are
bounded Boolean conjunctions.
\item \texttt{PolyCode}: a list of coefficient/exponent-vector pairs;
\texttt{eval}, \texttt{partial}, \texttt{jacobian},
\texttt{matVec}, and \texttt{matMul} are explicit PR folds.  A
\texttt{GoodPoly} predicate checks every arity and carrier condition.
\item \texttt{liftState P U a r}: the recursion (L), with a default value
on malformed input; \texttt{LiftHyp} is the bounded conjunction (H).
\item \texttt{JetCode N} and \texttt{ModCode q N}: the concrete
implementations (4) and (5), each with proofs of the abstract interface.
\texttt{HurwitzCode N} replaces ordinary convolution by (6).
\end{enumerate}

The main mathematical targets are the following explicit propositions,
with all variables understood as numerical codes:
\begin{align*}
 &\operatorname{LiftHyp}(N,P,U,a)\ \longrightarrow\
   \operatorname{Eval}(P,\operatorname{liftState}(P,U,a,N-1))=0,\\
 &\operatorname{LiftHyp}(N,P,U,a)\wedge P(x)=0\wedge x-a\in F^1
      \ \longrightarrow\
   x=\operatorname{liftState}(P,U,a,N-1).
\end{align*}
Their deep endpoints are terms of
\texttt{Prf PRAAx [] (allClose LiftExistence)} and
\texttt{Prf PRAAx [] (allClose LiftUniqueness)}, after the supplied
presentation laws have been discharged or retained explicitly as premises.
They are not statements that an arbitrary host ring can be automatically
turned into an arithmetic presentation.

The required intermediate lemmas are, in order: finite sum and matrix
filtration closure; the uniform finite Taylor expansion in
Lemma~\ref{lem:filtered-taylor}; residual precision improvement (1);
the Boolean invariant (3); difference improvement (2); homomorphism
commutation with finite evaluation; and the two concrete implementations.
For Hurwitz equations add Pascal's identity, the triple-binomial identity,
coefficient causality for \texttt{eval}, and the bounded-prefix invariant
for (9).  Each item has its numerical definition and its exact induction
formula specified above.  The general theorem and these implementations
are supplied here as human proofs and a detailed formalization plan;
they are not reported as already compiled deep iPRA derivations.
